\section{Introduction}
\label{sec:introduction}

Power-flow feasibility asks whether the electrical laws and operating limits
of a network can hold simultaneously. Even without switching decisions,
power is quadratic in voltage, so feasibility is a nonconvex polynomial
problem. Two questions follow: how hard is it to decide feasibility exactly,
and how much does a small numerical residual establish? We answer both for a
resistive model with rational conductances, positive voltage intervals, and
signed injection intervals, and transfer the exact classification to several
precisely specified AC models. The results identify worst-case obstructions
to exact certification; they make no claim about the typical difficulty of
operating networks.

For AC networks, exact feasibility is known to be NP-hard
\citep{BienstockVerma2019,LehmannEtAl2016}.
NP-hardness does not settle whether every feasible instance has a short
certificate. The relevant class is $\ER$, the problems that reduce in
polynomial time to the existential theory of the reals. It satisfies
$\NP\subseteq\ER\subseteq\PSPACE$, and whether $\NP=\ER$ is open. The
classical $\ER$-hardness results go back to Mn\"ev's universality theorem
\citep{Mnev1988} and Shor's reduction for stretchability of pseudolines
\citep{Shor1991}; later reductions from bounded arithmetic, as for the art
gallery problem \citep{AbrahamsenEtAl2018,AbrahamsenEtAl2022}, extend
$\ER$-hardness to many geometric problems. \citet{SchaeferEtAl2024} survey this literature and
the relation of $\ER$ to exact real-valued certificates. To the best of
our knowledge, the prior power-flow results discussed below do not give an
$\ER$-completeness theorem for the restricted resistive model studied here,
or an electrical realization of bounded arithmetic that preserves the
entire solution set.

\subsection{Main results}

\begin{enumerate}
\item \emph{$\ER$-completeness under simultaneous restrictions.}
Exact resistive feasibility is $\ER$-complete (Theorem~\ref{thm:rpf-complete}).
For each fixed integer $r$, hardness holds simultaneously on connected
simple planar bipartite graphs of maximum degree three and girth at least
$r$, with unit conductance and interval endpoints from fixed finite rational
sets depending only on $r$ (Theorem~\ref{thm:structural}). Thus neither
large numerical coefficients nor dense graphs are required. A polynomial-size certificate for every feasible instance, verifiable in
deterministic polynomial time, would imply $\NP=\ER$
(Corollary~\ref{cor:rpf-np}).

\item \emph{Universality and algebraic degree.}
The reduction preserves the entire solution set by rational maps in both
directions (Proposition~\ref{prop:rational-bijection}). Combined with a
bounded-arithmetic lemma (Lemma~\ref{lem:basic-arithmetic}) and an invariance
argument (Proposition~\ref{prop:basic-invariance}), this shows that a
compact semialgebraic set defined over $\Q$ is rationally equivalent to the
feasible voltage set of a rational network if and only if it is basic
closed over $\Q$, that is, described by a conjunction of polynomial
equalities and weak inequalities with rational coefficients
(Theorem~\ref{thm:universality}). Via semialgebraic triangulation, every compact
semialgebraic set occurs up to semialgebraic homeomorphism
(Theorem~\ref{thm:topological-universality}).
A rational network can have a unique feasible voltage vector with a
designated voltage generating any prescribed real algebraic number field
(Theorem~\ref{thm:algebraic-degree}). All three results hold under the
restrictions of Theorem~\ref{thm:structural}.

\item \emph{Residuals and certificates.}
For the ordinary construction, network and source residuals agree up to a
factor linear in the number of equations
(Theorem~\ref{thm:residual-transfer}). An explicit infeasible family with
$O(k)$ buses, degree at most three, and fixed numerical data admits
injection residuals below $2^{-2^k}$ while respecting every voltage bound
exactly (Theorem~\ref{thm:tiny-residual}). A general polynomial-minimum
bound shows that this doubly exponential scale is the worst case for the
fixed-data bounded-degree class (Proposition~\ref{prop:residual-separation}).
The residual family is not asserted to have the additional restrictions of
the planar construction. A universal residual threshold for exact
feasibility can therefore require exponentially many accuracy bits in the
number of buses. This limits certification by residual size alone; it is
not a running-time lower bound for exact algorithms. Under an explicit gap
promise, short rational certificates exist
(Proposition~\ref{prop:approx-certificates}).

\item \emph{AC angle semantics and transfers.}
A local cosine inequality constrains the principal angle of a line but need
not admit real bus angles whose ordinary differences satisfy the same limit:
a cycle can wind around the origin while every local inequality holds. We
encode the zero-winding requirement by a quadratic existential formula with
$O(|N|+|E|)$ variables, predicates, and monomial occurrences, allowing
unequal voltage magnitudes and every rational cosine limit in $(-1,1]$
(Theorem~\ref{thm:ac-membership}). For resistive lines, an energy identity
forces equal real angles when all reactive injections vanish and all line
differences have magnitude less than $\pi$ (Lemma~\ref{lem:equal-angles}).
This transfers $\ER$-completeness to AC feasibility with real-angle limits
(Theorem~\ref{thm:ac-complete}), including one-sided reactive intervals of
positive width (Corollary~\ref{cor:one-sided-reactive}). Separate transfers
cover reference-fixed bus-angle boxes (Corollary~\ref{cor:box-hardness})
and strictly positive principal-angle windows that shrink with the number
of buses (Corollary~\ref{cor:principal-hardness}); the latter use
polynomial-bit rational cosines, which need not come from a fixed finite
set. The graph restrictions of Theorem~\ref{thm:structural} carry over
(Corollary~\ref{cor:structural-ac}). Principal angles alone admit winding
solutions with no resistive counterpart
(Proposition~\ref{prop:winding-counterexample}). An energy estimate also
bounds the angle deviation and active-power error caused by small reactive
injections (Proposition~\ref{prop:reactive-stability}).
\end{enumerate}

\paragraph{Key idea.}
A bus with prescribed voltage and injection imposes an affine relation on
neighboring voltages. Alternating complement paths supply distinct copies
of a source variable without increasing the degree beyond three, and a
small nonlinear gadget enforces inversion. Every solution of the bounded
arithmetic source problem has exactly one network extension, and designated
bus voltages recover the source variables (Section~\ref{sec:resistive}).
A bounded arithmetic crossover gives planarity, free-injection connectors
give connectedness, and harmonic subdivisions give unit conductance,
bipartiteness, and prescribed girth (Section~\ref{sec:structural}).

\subsection{Relation to prior work}

\paragraph{Complexity of power flow.}
\citet{BienstockVerma2019} prove strong NP-hardness in a lossless AC model
with fixed voltage magnitudes and unconstrained reactive power.
\citet{LehmannEtAl2016} prove NP-hardness on trees, already using star
networks, in a fixed-magnitude model with active and reactive constraints.
Our classification identifies the full existential-real complexity of a
different subclass: resistive lines and variable magnitudes, with reactive
constraints that force equal angles under the stated angle convention.
It neither strengthens the tree restriction in their theorem nor proves
$\ER$-hardness for the lossless fixed-magnitude model.
Table~\ref{tab:prior-complexity} compares the models.

\begin{table}[t]
\centering
\small
\caption{Exact feasibility classifications for different power-flow models.
The restrictions within each row hold simultaneously.}
\label{tab:prior-complexity}
\begin{tabular}{@{}p{0.23\linewidth}p{0.49\linewidth}p{0.22\linewidth}@{}}
\toprule
Result & Model and restrictions & Classification \\
\midrule
\citet{BienstockVerma2019} & Lossless AC; fixed voltage magnitudes;
unconstrained reactive power & Strongly NP-hard \\
\addlinespace
\citet{LehmannEtAl2016} & AC on stars; fixed voltage magnitudes;
active and reactive constraints & NP-hard \\
\addlinespace
Theorem~\ref{thm:structural} & Resistive; connected, simple, planar,
bipartite; degree at most three; any fixed girth bound; unit conductance;
fixed finite endpoint sets & $\ER$-complete \\
\addlinespace
Theorem~\ref{thm:ac-complete} & Resistive-line AC; variable magnitudes;
zero reactive injections; real line differences in $[-\pi/2,\pi/2]$
& $\ER$-complete \\
\bottomrule
\end{tabular}
\end{table}

\paragraph{Resistive models and relaxations.}
Semidefinite \citep{LavaeiLow2012} and second-order cone
\citep{FarivarLow2013,GanLow2014} relaxations of optimal power flow are
exact under sufficient conditions, which our hard instances need not
satisfy. The resistive equations are the physical direct-current power-flow
model studied by \citet{GanLow2014}, not the linear ``DC approximation'' of
AC power flow. Their exactness conditions for resistive optimal power flow
include conditions involving voltage upper bounds and injection lower
bounds. Our model allows independent voltage and injection intervals, both of which may
be singletons at the same bus, and the reduction uses buses with both
prescribed voltage and prescribed injection. Section~\ref{sec:foundations}
states these input choices; we do not establish hardness after removing
the prescribed quantities used by the construction.

A related resistive model fixes source voltages and asks which
constant-power demands can be supplied at load buses.
\citet[Theorems 3.18 and 3.22]{JeeningaEtAl2023} characterize exact
feasibility and its interior by matrix-inequality alternatives, and
establish convexity of the feasible demand set. Their model has no
independent operating voltage bounds or injection restrictions at the
fixed-voltage sources. Those restrictions distinguish our input model;
signed demands alone do not. An exact matrix-inequality characterization
also does not, by itself, give a polynomial-time Turing algorithm for exact
feasibility.

\paragraph{Angle conventions and winding.}
Cycle consistency and winding are established features of power-flow
models. \citet[Theorem 2]{FarivarLow2013} characterize branch-flow angle
recovery by cycle sums that vanish modulo $2\pi$.
\citet{DelabaysEtAl2016} relate loop currents to winding numbers, and
\citet{JafarpourEtAl2022} develop winding cells and uniqueness, modulo a
common rotation, within each cell for monotone flow networks on the torus.
In our real-angle model, the chosen short line differences must instead sum
to zero as real numbers, which selects the zero-winding cell. We do not
claim that criterion, the existence of winding solutions, or uniqueness in
the monotone angular regime as new. Our additional step is to express the
required branch choices and zero winding in polynomial size with rational
coefficients, without taking an inverse trigonometric function as an exact
arithmetic primitive. This makes the $\ER$ membership statement precise.

The resistive-AC connection with zero reactive injections already appears
in \citet[Appendix B, Case 2]{LavaeiLow2012}; weighted sine-coupling
equilibria also arise in \citet{DorflerEtAl2013}. We give a self-contained
energy proof under the real-angle hypothesis needed here.

\paragraph{Universality and algebraic structure.}
\citet{MarecekEtAl2016} study polynomial formulations, generic complex
solution counts, and positive-dimensional solution sets of power flow. Our
universality results instead realize prescribed compact feasible sets, with
operating inequalities, in the restricted resistive family; they do not
assert that such behavior is generic. The bounded arithmetic source problem
is due to \citet{AbrahamsenEtAl2018,AbrahamsenEtAl2022}, and the planar
crossover is from \citet{DobbinsEtAl2023}. For rational universality we use
the conjunction-only arithmetic ideas of \citet{AbrahamsenMiltzow2019}, with
explicit range checks in Appendix~\ref{app:arithmetic}. The broader
rational-equivalence statement in Theorem~1 of that preprint includes
arbitrary compact semialgebraic sets, and its Definition~4 uses the same
rational-coordinate notion as here: the forward and inverse maps must be
defined on their respective sets.
Proposition~\ref{prop:basic-invariance} and the three-quadrant example rule
out that broader scope under this definition. We therefore prove the
basic-closed version directly, without its Boolean preprocessing.
Semialgebraic triangulation \citep{OhmotoShiota2017} gives the separate
topological conclusion. The new application is the electrical realization
under all the stated restrictions, not universality of bounded arithmetic
or triangulation itself.

\paragraph{Residuals and approximate certificates.}
Doubly exponential separation is known for general polynomial systems:
\citet{JeronimoEtAl2013} give both a polynomial-minimum bound and a
repeated-power example showing the need for that scale. Our contribution is
an explicit infeasible electrical family using only fixed data and bounded
degree, together with a quantitative bound for the actual arithmetic
gadgets. \citet{BienstockDelPiaHildebrand2023} study short rational points
and approximate certificates in polynomial optimization. Our gap-promise
certificate follows from a direct Lipschitz and rounding argument for the
bounded quadratic network model; we do not present that general
approximation principle as new.
\citet[Theorem 7 and Corollary 8 of the arXiv version]{BienstockMunoz2018}
give LP approximation schemes for polynomial network optimization under
bounded treewidth and bounded numbers of local variables and constraints,
including AC optimal power flow, with scaled feasibility and optimality
tolerances. Those guarantees do not decide exact feasibility.

\paragraph{Organization.}
Sections~\ref{sec:foundations}--\ref{sec:resistive} give the model and the
resistive reduction, Section~\ref{sec:ac} the AC encoding and transfers,
Section~\ref{sec:structural} the graph restrictions, and
Sections~\ref{sec:algebraic}--\ref{sec:numerical} the universality and
quantitative results. Appendices~\ref{app:arithmetic}
and~\ref{app:verification} contain the bounded arithmetic lemma and the
exact-arithmetic checks, which support, and do not replace, the proofs.
